\uuid{Lx3T}
\titre{Expérience de Michelson : tests sur la vitesse de la lumière}
\niveau{L3}
\module{Probabilité et statistique}
\chapitre{Tests sur la moyenne, test du chi2}
\sousChapitre{Tests d'hypothèses, intervalle de confiance}
\theme{Test de Student, intervalle de confiance, test de la variance, chi-deux, puissance, p-valeur, approximation de Fisher}
\auteur{Erwan l'Haridon}
\datecreate{2026-05-19}
\organisation{AMSCC}
\difficulte{3}

\contenu{
	
	\texte{
		En 1879, Michelson effectue 100 mesures de la vitesse de la lumière (en km$\cdot$s$^{-1}$). On modélise ces mesures comme des variables aléatoires i.i.d. gaussiennes d'espérance $\mu$ et d'écart-type $\sigma$ inconnus. Sur cet échantillon, on dispose des observations suivantes :
		$$
		\bar{x} = 299\,852 \text{ km$\cdot$s$^{-1}$}, \qquad s = 79 \text{ km$\cdot$s$^{-1}$} \quad (n = 100).
		$$
		L'estimateur $S$ utilisé ici est $S = \sqrt{\frac{1}{n-1}\sum_{i=1}^n (X_i - \overline{X})^2}$.
		
		À l'époque, la valeur communément admise était $c_{1879} = 299\,840$ km$\cdot$s$^{-1}$. Aujourd'hui, elle est établie à $c_{2026} = 299\,792$ km$\cdot$s$^{-1}$.
		
		Dans tout l'exercice, le risque de première espèce vaut $\alpha = 5\%$.
	}
	
	\begin{enumerate}
		
		\item
		\question{Tester l'hypothèse nulle que les mesures de Michelson sont conformes à la valeur $c_{1879} = 299\,840$ km$\cdot$s$^{-1}$ communément admise en 1879, puis tester si elles sont en accord avec la valeur établie aujourd'hui $c_{2026} = 299\,792$ km$\cdot$s$^{-1}$. Pour chaque test, construire la statistique, la région de rejet et conclure.}
		\indication{La variance étant inconnue, on utilise la statistique de Student ; pour $n = 100$, la table de Student n'étant pas disponible pour de grands échantillons, on approchera la loi de Student par une loi normale centrée réduite.}
		\reponse{
			\textbf{Test 1 : conformité avec $c_{1879}$.}
			
			On teste $H_0 : \mu = c_{1879}$ contre $H_1 : \mu \neq c_{1879}$. C'est un test bilatéral sur la moyenne d'un échantillon gaussien de variance inconnue. La statistique de test est :
			\[
			T = \frac{\bar{X} - c_{1879}}{S/\sqrt{n}} \sim \mathcal{S}t(n-1) = \mathcal{S}t(99) \quad \text{sous } H_0.
			\]
			Pour $n = 100$, on approche $\mathcal{S}t(99)$ par $\mathcal{N}(0,1)$. Rappelons que le \textbf{quantile d'ordre $p$} d'une loi est la valeur $q_p$ telle que $P(X \le q_p) = p$. La région de rejet au seuil de $5\%$ est :
			\[
			\mathcal{R} = \left]-\infty\,;\,-q_{0{,}975}\right[ \cup \left]q_{0{,}975}\,;\,+\infty\right[, \quad \text{avec } q_{0{,}975} \approx 1{,}96.
			\]
			Application numérique :
			\[
			t_{\text{obs}} = \frac{299\,852 - 299\,840}{79/\sqrt{100}} \approx 1{,}52.
			\]
			Comme $|t_{\text{obs}}| = 1{,}52 < 1{,}96$, on ne rejette pas $H_0$ au seuil de $5\%$ : les mesures de Michelson sont jugées compatibles avec $c_{1879}$.
			
			La $p$-valeur est :
			\[
			p\text{-val} = P_{H_0}(|T| > 1{,}52) \approx 2\Phi(-1{,}52) \approx 0{,}13 > 0{,}05.
			\]
			On confirme le non-rejet de $H_0$.
			
			Le risque de seconde espèce ne peut pas être calculé car l'hypothèse alternative $H_1 : \mu \neq c_{1879}$ est composite (elle ne spécifie pas une valeur exacte pour $\mu$).
			
			\medskip
			\textbf{Test 2 : conformité avec $c_{2026}$.}
			
			On teste $H_0 : \mu = c_{2026} = 299\,792$ contre $H_1 : \mu \neq c_{2026}$. Sous $H_0$, la même statistique de test suit approximativement une loi $\mathcal{N}(0,1)$, avec la même région de rejet.
			
			Application numérique :
			\[
			t_{\text{obs}} = \frac{299\,852 - 299\,792}{79/\sqrt{100}} \approx 7{,}59.
			\]
			Comme $|t_{\text{obs}}| = 7{,}59 \gg 1{,}96$, on rejette $H_0$. Les mesures de Michelson ne sont pas en accord avec la valeur moderne établie en 2026.
		}
		

\texte{On souhaite tester l'écart-type $\sigma$ de la mesure. Sous $H_0 : \sigma = \sigma_0$, la statistique de test usuelle est :
	$$
	K = (n-1)\,\frac{S^2}{\sigma_0^2} \sim \chi^2(n-1).
	$$
	La table du $\chi^2$ n'étant pas disponible pour $d = n-1 = 99$ degrés de liberté, on utilise l'approximation de Fisher : si $K \sim \chi^2(d)$, alors pour $d$ grand, la statistique $Z = \sqrt{2K} - \sqrt{2d-1}$ suit approximativement une loi normale centrée réduite $\mathcal{N}(0, 1)$. 
	
	On rappelle que le quantile d'ordre $p$ d'une loi est la valeur $q_p$ telle que $P(X \le q_p) = p$.}


	\item
	\question{À l'aide de cette approximation et d'une table de la loi normale centrée réduite, déterminer une approximation du quantile d'ordre $95\%$ de la loi $\chi^2(99)$, noté $q_{\chi^2(99),\, 0{,}95}$. }
	\reponse{
		Soit $K \sim \chi^2(99)$. On cherche la valeur $c$ telle que $P(K \le c) = 0{,}95$.
		En appliquant la transformation de Fisher avec $d = 99$ :
		\[
		P(K \le c) = P\!\left(\sqrt{2K} - \sqrt{197} \le \sqrt{2c} - \sqrt{197}\right) \approx 0{,}95.
		\]
		La variable $Z = \sqrt{2K} - \sqrt{197}$ étant approximativement distribuée selon une loi $\mathcal{N}(0, 1)$, on identifie le seuil à l'aide du quantile d'ordre $0{,}95$ de la loi normale :
		\[
		\sqrt{2c} - \sqrt{197} \approx z_{0{,}95} = 1{,}645.
		\]
		On en déduit :
		\[
		\sqrt{2c} \approx \sqrt{197} + 1{,}645 \approx 14{,}036 + 1{,}645 = 15{,}681.
		\]
		En élevant au carré :
		\[
		2c \approx (15{,}681)^2 \approx 245{,}90 \implies c \approx 122{,}95.
		\]
		L'approximation de Fisher donne $q_{\chi^2(99),\, 0{,}95} \approx 122{,}95$ (la valeur exacte est $123{,}22$).
	}
	
	\item
	\question{Tester $H_0 : \sigma = 90$ contre $H_1 : \sigma = 91$ au risque $5\%$. Calculer le risque de seconde espèce $\beta$.}
	\reponse{
		On effectue un test unilatéral à droite. Sous $H_0$, la statistique $Z = \sqrt{2K} - \sqrt{197}$ suit approximativement une loi $\mathcal{N}(0, 1)$. La région de rejet au seuil $5\%$ est :
		\[
		\mathcal{R} = \left\{ z_{\text{obs}} > 1{,}645 \right\}.
		\]
		Calcul de la statistique observée :
		\[
		k_{\text{obs}} = 99 \times \frac{79^2}{90^2} \approx 76{,}28 \implies z_{\text{obs}} = \sqrt{2 \times 76{,}28} - \sqrt{197} \approx 12{,}35 - 14{,}04 = -1{,}69.
		\]
		Comme $z_{\text{obs}} \approx -1{,}69 \le 1{,}645$, on ne rejette pas $H_0$.
		
		\textbf{Calcul du risque de seconde espèce $\beta$ :}
		
		On rejette $H_0$ si $K > q_{\chi^2(99),\, 0{,}95} \approx 122{,}95$ (d'après la question précédente), ce qui correspond à :
		\[
		S^2 > \frac{122{,}95 \times 90^2}{99} \approx 10\,059{,}55 \text{ km}^2\cdot\text{s}^{-2}.
		\]
		Sous $H_1 : \sigma = 91$, la variable $K' = 99\,\dfrac{S^2}{91^2}$ suit une loi $\chi^2(99)$. Par l'approximation de Fisher, $Z' = \sqrt{2K'} - \sqrt{197} \sim \mathcal{N}(0, 1)$.
		
		Le risque de seconde espèce $\beta$ est la probabilité de ne pas rejeter $H_0$ alors que $H_1$ est vraie :
		\[
		\beta = P_{H_1}\!\left(S^2 \le 10\,059{,}55\right) = P_{H_1}\!\left(K' \le 99 \times \frac{10\,059{,}55}{91^2}\right) = P_{H_1}\!\left(K' \le 120{,}26\right).
		\]
		En appliquant la transformation de Fisher sous $H_1$ :
		\[
		\beta \approx P\!\left(Z' \le \sqrt{2 \times 120{,}26} - \sqrt{197}\right) = P\!\left(Z' \le 15{,}51 - 14{,}04\right) = P\!\left(Z' \le 1{,}47\right).
		\]
		À l'aide de la table de la loi normale centrée réduite :
		\[
		\beta \approx \Phi(1{,}47) \approx 0{,}93.
		\]
		La puissance du test est faible ($1 - \beta \approx 7\%$) : l'échantillon ne permet pas de distinguer de manière fiable $\sigma = 90$ de $\sigma = 91$.
	}
	
\end{enumerate}
		
}